\subsection{Model I (G/Gauge A)}

\subsubsection{Identification}
\paragraph{Kinematics}
Choose the rigid/axis fields so that all \(x\)-polynomial, \(\bm{r}\)-affine parts of Saint–Venant are carried by \(\bm u=\bm x(x)-x\,\FrameBasis\) and \(\bm\theta\):

\begin{equation}
\boxed{
\begin{aligned}
\bm u(x)&=
x\,a_\varepsilon\,\FrameBasis
-\tfrac{1}{2} x^2\,\bm a_\Sigma\;-\;\tfrac16 x^3\,\bm{b}_{\Section}
\;-\;\underbrace{\left[\tfrac{1}{2} x^2\,(\FrameBasis\otimes\CentroidLocation)\bm{b}_{\Section}\right]}_{\text{uniform}},\\[2pt]
\bm\theta(x)&= x\,(a_\vartheta-c_\vartheta)\,\FrameBasis- \,x\,(\FrameBasis\times\bm a_\Sigma)\;-\;\tfrac{1}{2} x^2\,(\FrameBasis\times\bm{b}_{\Section}).
\end{aligned}}
\end{equation}
The uniform axial term in brackets is optional and may be absorbed into \(\bm u\) to avoid introducing a separate uniform-axial warping shape.

\paragraph{Deformations}
The deformations \(\bm{\gamma}=\bm u'+\FrameBasis\times\bm\theta\) and \(\bm{\kappa}=\bm\theta'\) are:
\[
\boxed{
\begin{aligned}
\bm{\gamma}(x)&=\ a_\varepsilon\,\FrameBasis,\\[2pt]
\bm{\kappa}(x)&=
(a_\vartheta-c_\vartheta)\,\FrameBasis
\ -\,(\FrameBasis\times\bm a_\Sigma)\;-\;x\,(\FrameBasis\times\bm{b}_{\Section})
\end{aligned}}
\]
Hence the alignment shear deformation \(\bm{\gamma}_\perp\equiv 0\) vanishes. 
All \(\bm b_\Sigma\)-driven shear-like content is realized via warping amplitudes.
The flexural curvature driven by \(\bm{b}_{\Section}\) is linear in \(x\), as in SV.

\paragraph{Warping Basis}
Use the canonical \(x\)-independent SV shapes:
\[
\bm\Phi=\big[
\bm{W}^{(a)}\mathbf{e}_y,\ \bm{W}^{(a)}\mathbf{e}_z,\ \bm{W}^{(\varepsilon)},\
\bm{W}^{(b)}\mathbf{e}_y,\ \bm{W}^{(b)}\mathbf{e}_z,\
\FrameBasis\,\psi_{y},\ \FrameBasis\,\psi_{z},\ \FrameBasis\,\varphi,\ 
[\;\FrameBasis\;]
\big],
\]
representing the 9 canonical bases defined by \eqref{eq:def-eight-shapes}.

\paragraph{Amplitudes}
\[
\boxed{
\begin{aligned}
&\bm{\alpha}_{(a)}\mathbf{e}_y=a_y,\\
&\bm{\alpha}_{(a)}\mathbf{e}_z=a_z,\\
&\bm{\alpha}_{(\varepsilon)}=a_\varepsilon,\\[2pt]
&\begin{bmatrix}\alpha(\bm{W}^{(b)}\mathbf{e}_y)\\ \alpha(\bm{W}^{(b)}\mathbf{e}_z)\end{bmatrix}=\bm{\eta}(x),\\
&\begin{bmatrix}\alpha(\FrameBasis\,\psi_y)\\ \alpha(\FrameBasis\,\psi_z)\end{bmatrix}=\bm{\alpha}_{\psi}(x),\\
&\alpha_\varphi=a_\vartheta+c_\vartheta.
\end{aligned}}
\]
\(\bm{\alpha}_{\psi},\bm{\eta}\) are \emph{independent} internal fields here; for the SV field below they will specialize to \(\bm{\beta}=\bm{b}_{\Section},\ \bm{\eta}=x\,\bm{b}_{\Section}\). 
If the bracketed rigid term is omitted from \(\bm{u}\) one additionally has 
\[
\alpha(\FrameBasis)=-\tfrac{1}{2}\,x^2\,(\CentroidLocation\cdot\bm{b}_{\Section}).
\]

\begin{equation*}
\boxed{
\begin{aligned}
&\bm{\alpha}_{(\varepsilon)}~~=a_\varepsilon,\\
&\bm{\alpha}_{(a)}=~~\bm{a}_{\Section},\\
&\bm{\alpha}_{(b)}=x\,\bm{b}_{\Section},\\
&\bm{\alpha}_{\WarpShear}=\bm{b}_{\Section} &\in\mathbb{R}^2,\\
&\bm{\alpha}_\varphi=a_\vartheta+c_\vartheta.\\
&\bm{\alpha}_{\mathbf{1}}(x)=-\tfrac{1}{2} x^2\,(\CentroidLocation\cdot\bm{b}_{\Section}),\\
\end{aligned}}
\end{equation*}


\paragraph{Remarks}
\begin{itemize}
\item The kinematic choice \(\bm{\gamma}_\perp\equiv \mathbf{0}\) avoids introducing linear-axial shapes \(y\,\FrameBasis,z\,\FrameBasis\) and keeps the \(\psi\)-pair amplitudes equal to \((b_y,b_z)\) exactly.  
\item No uniqueness is asserted; other gauges (e.g., introducing a constant shear offset in \(\bm\theta\)) can redistribute content between alignment and warping fields without changing the final SV field.
\item With \(\bm\Phi\) orthogonal to rigid modes and gauges for \(\varphi,\psi\) fixed, the coefficients multiplying the independent rigid patterns
\(\{\FrameBasis\times\bm{r},\ \FrameBasis\otimes\bm{r},\ \text{constant in }\bm{r}\}\)
are uniquely determined, yielding the boxed \(\bm\theta(x)\) and \(\bm u(x)\).
\end{itemize}

\subsubsection{Constitution}
\emph{Primary unknowns.}\;
\(\bm x(x)\) and \(\bm\theta(x)\), plus internal warping amplitudes
\(\bm{\alpha_{\psi}}(x)\in\mathbb{R}^2\) (flexural axial \(\psi\)-pair) and
\(\bm{\eta}(x)\in\mathbb{R}^2\) (bending–Poisson pair).

\emph{Generalized strains.}\;
\(\epsilon=\FrameBasis\cdot\bm u'=u_x'\),\;
\(\bm{\kappa}=\bm\theta'\),\;
\(\varkappa=\FrameBasis\cdot\bm\theta'\).

\paragraph{Material Stress/Strain}
The shear stress field induced by amplitudes $\bm{\alpha}_{\psi}, \bm\eta \in \mathbb{R}^2$ is
$$
\bm\tau(r)=G\left((\mathbf{1} + \nabla \bm{\psi}) \bm\alpha_{\psi}
+(\mathbf{1} + \nabla \bm{\psi})\bm{W}^{(b)}\bm\eta
\right),
$$
\paragraph{Resultants}
\[
V=\int \bm{\tau}\,\mathrm{d}A = \underbrace{\int_{\Sigma} G E_\psi \, d a }_{K_\psi}\,\bm\alpha_{\psi} 
+\underbrace{\int_{\Sigma} G E_b \,d a}_{K_b}\,\bm\eta
\]
\emph{1D energy density.}\; By inserting the kinematics in 3D linear elasticity and integrating over \(\Sigma\):
\[
\mathcal{U}
=\tfrac{1}{2} EA\,\epsilon^2
+\tfrac{1}{2} \bm{\kappa}\cdot E\mathbf{I}\,\bm{\kappa}
+\tfrac{1}{2} GJ\,\varkappa^2
+\tfrac{1}{2} \bm\alpha_{\psi} \cdot \mathbf{K}_\psi\,\bm\alpha_{\psi}
+\tfrac{1}{2} \bm\eta \cdot \mathbf{K}_b\,\bm\eta
+\text{(known small couplings, e.g.\ from }c_\vartheta).
\]
Here \(\mathbf{K}_\psi,\mathbf{K}_b\in\mathbb{R}^{2\times 2}\) are symmetric positive-definite matrices obtained from SV warping strains:
their entries are cross-section integrals of the \(\psi\)- and \(\bm{W}^{(b)}\)-induced strain–stress products.

\emph{Resultants and constitutive blocks.}\;
\[
N=EA\,\epsilon,\qquad
\bm{M}=\mathbf{EI}\,\bm{\kappa},\qquad
T=GJ\,\varkappa,
\]
and the shear force vector \(\bm V\in\mathbb{R}^2\) emerges from the warping modes as
\[
\bm V=\mathbf{K}_\psi\,\bm\alpha_{\psi}\;+\;\mathbf{K}_b\,\bm\eta.
\]
Linear superposition follows from linear elasticity; the matrices are the same as in the energy.

\emph{Field equations (static).}\;
With distributed loads \((n_x,\bm q,\tau)\) and moments \(\bm m\),
\[
\begin{aligned}
N'+n_x=0,\\
\bm{M}'+\FrameBasis\times N\,\FrameBasis+\bm m - \bm V=\mathbf{0},\\
T'+\tau=0,\\
\bm V'+\bm q=\mathbf{0}.
\end{aligned}
\]
The internal amplitudes are \emph{algebraic}:
\[
\frac{\partial \mathcal{U}}{\partial \bm\alpha_{\psi}}=\mathbf{K}_\psi\,\bm\alpha_{\psi}-\bm V=\mathbf{0},\qquad
\frac{\partial \mathcal{U}}{\partial \eta}=\mathbf{K}_b\,\bm\eta-\mathbf{0}=\mathbf{0},
\]
so, in the absence of additional driving for \(\bm{\eta}\), the optimal choice is \(\bm{\eta}=\mathbf{0}\) and \(\bm\alpha_{\psi}=\mathbf{K}_\psi^{-1}\bm V\).
If one wishes to reproduce the exact SV tip-load polynomial \(-x\,\nu(\cdots)\bm{b}_{\Section}\) in the kinematics, \(\bm{\eta}\) can be retained as an internal field and determined by a weak constraint or calibration so that under a constant shear resultant \(\bm V\) (cantilever), \(\bm{\eta}(x)=x\,\mathbf{K}_b^{-1}\bm V\).
Either choice leaves the \emph{1D} problem local and algebraic in \((\bm\alpha_{\psi},\bm\eta)\).

\emph{Elimination (single-field form).}\;
Eliminating \(\bm{\alpha_{\psi}}\) and \(\bm{\eta}\) yields a shear-energy term in the 1D functional:
\[
\mathcal{U}_\text{shear}=\tfrac{1}{2} \bm V \cdot \mathbf{K}_\psi^{-1}\bm V
\quad(\text{or } \tfrac{1}{2} \bm V^\top \mathbf{K}_\text{eff}^{-1}\bm V\ \text{if }\eta\text{ is retained}),
\]
with \(\mathbf{K}_\text{eff}^{-1}\) assembled from \(\mathbf{K}_\psi^{-1}\) and \(\mathbf{K}_b^{-1}\).
No shear-correction factor is needed; shear stiffnesses are exact for \(\Sigma\).

\emph{Recovery of the SV cantilever.}\;
For a tip shear \(\bm V\equiv \mathbf{V}_0\) (cantilever),
\(\bm\alpha_{\psi}=\mathbf{K}_\psi^{-1}\mathbf{V}_0\) (constant), and, if retained, \(\bm{\eta}(x)=x\,\mathbf{K}_b^{-1}\mathbf{V}_0\).
Substituting in the kinematics reproduces the SV polynomials:
constant \(\psi\)-amplitude for the axial warping and linear-in-\(x\) \(\bm{W}^{(b)}\)-amplitude for the in-plane Poisson warping,
while the alignment produces the flexure/torsion/extension blocks exactly.

\emph{Remarks.}\;
\begin{itemize}
\item Shear is represented through internal SV modes with section-exact stiffness matrices, avoiding ad hoc shear-correction factors.
\item The formulation remains strictly 1D and local; \(\bm\alpha_{\psi},\bm\eta\) are algebraic and can be kept or eliminated without affecting equilibrium structure.
\end{itemize}

